\chapter*{LAMPIRAN}
\addcontentsline{toc}{chapter}{LAMPIRAN}
\appendix
\inappendixtrue

\chapter{DOKUMENTASI}

Berikut adalah beberapa dokumentasi tangkapan layar dari percobaan yang dilakukan pada sistem.

\insertimage{assets/contoh-gambar.png}{Tampilan awal basis data}{fig:lampiran-1}

\insertimage{assets/contoh-gambar.png}{Tabel audit setelah disisipkan trigger}{fig:lampiran-2}

\chapter{SOURCE CODE}

\begin{lstlisting}[language=SQL, caption={Skrip lengkap pembuatan tabel \texttt{transaksi\_audit}}]
CREATE TABLE transaksi_audit (
    id          INT AUTO_INCREMENT PRIMARY KEY,
    id_trans    INT NOT NULL,
    aksi        VARCHAR(10) NOT NULL,
    data_lama   TEXT,
    data_baru   TEXT,
    waktu       DATETIME DEFAULT NOW()
);
\end{lstlisting}

\chapter{DATA PENGUJIAN}

Contoh data pengujian tertera pada tabel berikut.

\begin{table}[H]
    \centering
    \caption{Sampel Data Masukan Uji Coba}
    \label{tab:data-lampiran}
    \begin{tabularx}{\textwidth}{@{} c c X @{}}
        \toprule
        \textbf{No.} & \textbf{Nama Kasus} & \textbf{Keterangan} \\
        \midrule
        1 & Tambah data baru & INSERT ke tabel transaksi \\
        2 & Ubah transaksi & UPDATE nilai nominal \\
        3 & Hapus transaksi & DELETE berdasarkan ID \\
        \bottomrule
    \end{tabularx}
\end{table}
